\documentclass[10pt,a4paper,openany,twoside, prof]{book}

\usepackage{layout_cours}
\begin{document}

\chapnumber{12} % numéro du chapitre
\titre{Equation de droite} % titre du chapitre
\classe{Seconde} % classe
% \entetegal

\subsection*{37 p.136}
\begin{enumerate}[label=\alph*)]
   \item $\vecu \coordvp{1}{\sqrt{2}}$
   \item $\vecu \coordvp 1 1$
   \item $\vecu \coordvp 2 4$
   \item $\vecu \coordvp {-6} {2}$
\end{enumerate}

\subsection*{44 p.136}
$\v{GH}\coordvp{-2.25}{-0.1}$

$\v{GM}\coordvp{x-3}{y-0.2}$

On résout :
\begin{align*}
   \v{GH} \text{ et } \v{GM} \text{ sont colinéaires}
   & \iff \det{\v{GH}}{\v{GM}} = 0 \\
   &\iff -2.25(y-0.2)+0.1(x-3) = 0 \\
   &\iff 0.1x-2.25y+0.15 = 0 \\
   &\iff 2x-45y+3 = 0 \\
\end{align*}

Donc la droite $\mathcal{D}$ admet les équation cartésiennes $0.1x-2.25y+0.15 = 0$ ou $2x-45y+3 = 0$.

\bigskip
\textbf{2ème méthode}
On sait que $\v{GH}$ est un vecteur directeur. On en déduit les coefficients $a=-0.1$ et $b=2.25$ de l'équation cartésienne $ax+by+c=0$.

De plus, puisque $G\in \mathcal{D}$, on peut déterminer le coefficient $c$.

On résout :

\begin{align*}
   -0.1\times 3 +2.25 \times 0.2 +c=0 \iff c=-0.15
\end{align*}

On en déduit l'équation cartésienne de $\mathcal{D}$, $-0.1x+2.25y-0.15=0$.

\subsection*{73 p.136}
On résout le système : 
\begin{align*}
   (S) 
   &\iff \syst{2x+1.5y &=& 4.25 \\ 2.5x+5y &=& 8.75 }
   &\iff \syst{20x+15y &=& 42.5 \\ 20x+40y &=& 70 }
   &\iff \syst{20x+15y &=& 42.5 \\ 25y &=& 27.5 } \\
   &\iff \syst{x+0.75y &=& 2.125 \\ y &=& 1.1 }
   &\iff \syst{x &=& 2.125-0.75\times 1.1 \\ y &=& 1.1 }
   &\iff \syst{x &=& 1.3 \\ y &=& 1.1 }
\end{align*}
\end{document}

